% CUMCM 电子版：第一页为摘要页，正文自第 2 页起，无目录，正文 <= 20 页；匿名。
\documentclass[12pt,a4paper]{ctexart}
\begin{document}

\begin{center}
{\Large\bfseries 薄膜厚度反演的脱敏回归稿}\\
{\small （赛题年份 2025；不得据此推断提交规则年份）}
\end{center}

\textbf{摘\quad 要}

本文针对薄膜测厚问题，首先建立干涉模型，随后给出反演步骤，最后讨论可靠性。主要结果用概括句写出，不按子问对应。

\vspace{0.4em}
\textbf{关键词}：薄膜；干涉；反演

\vspace{0.5em}
\textbf{Abstract}\quad
We first build a model, then invert the spectrum, and finally discuss reliability.

\newpage
\section{问题重述}
题目要求回答三个子问。

\section{问题分析}
分析从预处理到拟合的流程。

\section{模型评价}
\subsection{优点}
\begin{itemize}
\item 步骤清楚
\item 结果可复述
\item 结构完整
\end{itemize}
\subsection{局限}
\begin{itemize}
\item 未交代参数来源
\item 未给出验证细节
\item 未传播不确定度
\end{itemize}
\subsection{改进方向}
\begin{itemize}
\item 补充文献引用
\item 补充敏感性
\item 补充残差
\end{itemize}

关键公式强调如下：
\boxed{d=1/(2 n \Delta\nu)}

\begin{thebibliography}{99}
\bibitem{param-source} Some Archive. Optical constants. \url{https://example.invalid/optics}
\bibitem{method-paper} A.~Author. A method paper. \textit{J.~Example} 2024.
\end{thebibliography}
\end{document}
